\section{计算机视觉简史}

本节按时间线梳理计算机视觉领域的关键里程碑，从神经科学的早期发现到深度学习革命。

\subsection{1950s--1960s：视觉研究的萌芽}

\subsubsection{Hubel \& Wiesel 的视觉通路研究（1959）}

1950年代，Hubel 和 Wiesel 对哺乳动物视觉通路进行了开创性的电生理实验。他们将电极插入麻醉猫的初级视觉皮层（primary visual cortex，位于后脑勺附近），发现了两个重要事实：

\begin{figure}[H]
\centering
\includegraphics[width=0.95\textwidth]{slides-images/slide-017.jpg}
\caption{Hubel \& Wiesel 的实验：电极记录猫初级视觉皮层神经元的感受野（receptive field），发现神经元响应特定方向的边缘\protect\footnotemark}
\end{figure}
\footnotetext{Slides 第18页（Part 1）。}

\begin{importantbox}{Hubel \& Wiesel 的两大发现}
\begin{enumerate}
    \item \textbf{感受野（Receptive Field）}：初级视觉皮层中的每个神经元只响应视野中一个\textbf{局部区域}的特定简单模式（如特定方向的边缘）。
    \item \textbf{层级结构（Hierarchical Processing）}：视觉通路是分层的——浅层神经元检测简单特征（方向边缘），深层神经元检测更复杂的特征（角点、物体部件），最终构成一个庞大的计算网络。
\end{enumerate}
这两个发现后来深刻影响了卷积神经网络的设计。Hubel 和 Wiesel 因此于 1981 年获得诺贝尔医学奖。
\end{importantbox}

\subsubsection{计算机视觉的诞生（1963--1966）}

\begin{itemize}
    \item \textbf{1963年}：Larry Roberts 完成了被公认为计算机视觉领域的第一篇博士论文，研究如何从图像中理解简单几何形状的表面、角点和特征。
    \item \textbf{1966年}：MIT 的 Seymour Papert 发起了著名的 \textbf{Summer Vision Project}，目标是``在一个暑假内解决计算机视觉问题''。
\end{itemize}

\begin{figure}[H]
\centering
\includegraphics[width=0.95\textwidth]{slides-images/slide-020.jpg}
\caption{MIT Summer Vision Project（1966）：Seymour Papert 试图在一个暑假内``解决视觉''，这反映了早期 AI 研究者的过度乐观\protect\footnotemark}
\end{figure}
\footnotetext{Slides 第21页（Part 1）。}

\begin{warningbox}{AI 的过度乐观传统}
就像 AI 历史上的许多故事一样，早期研究者严重低估了视觉问题的难度。事实上，直到今天——将近60年后——计算机视觉仍未被完全``解决''。这提醒我们对技术的能力保持谦逊。
\end{warningbox}

\subsection{1970s--1980s：理论框架与特征提取}

\subsubsection{David Marr 的视觉计算理论}

1970年代，David Marr 撰写了计算机视觉领域的奠基之作。他提出了一个系统化的视觉处理框架，受到神经科学和认知科学的深刻启发：

\begin{figure}[H]
\centering
\includegraphics[width=0.95\textwidth]{slides-images/slide-022.jpg}
\caption{David Marr 的视觉处理层级：从输入图像到 Primal Sketch（边缘等底层特征），到 2.5D Sketch（前景/背景深度分离），再到完整的 3D 表示\protect\footnotemark}
\end{figure}
\footnotetext{Slides 第23页（Part 1）。}

\begin{knowledgebox}{David Marr 的三级视觉表示}
\begin{itemize}
    \item \textbf{Primal Sketch}：提取边缘、纹理等底层特征（类似 Hubel \& Wiesel 发现的初级视觉皮层功能）
    \item \textbf{2.5D Sketch}：分离前景和背景，估计局部表面朝向和深度
    \item \textbf{3D Model}：恢复完整的三维几何表示——Marr 认为这是视觉理解的终极目标
\end{itemize}
David Marr 英年早逝，但他的理论框架深刻影响了后续几十年的研究方向。
\end{knowledgebox}

\subsubsection{物体识别与边缘检测的早期尝试}

1970年代至1980年代，研究者开始尝试用计算方法识别物体：

\begin{itemize}
    \item \textbf{广义圆柱体}（Generalized Cylinders，Brooks \& Binford，斯坦福）：用简单几何体组合来表示物体
    \item \textbf{组合模型}：用部件（parts）的组合来描述人体和其他物体
    \item \textbf{边缘检测}（Edge Detection）：从数字图像中提取边缘信息
\end{itemize}

\begin{figure}[H]
\centering
\includegraphics[width=0.95\textwidth]{slides-images/slide-024.jpg}
\caption{1970年代的物体识别方法：广义圆柱体表示和组合模型，用简单几何形状描述复杂物体\protect\footnotemark}
\end{figure}
\footnotetext{Slides 第25页（Part 1）。}

然而，这些方法的能力非常有限，仅能处理高度简化的场景，距离真正的``视觉理解''还有很大距离。这一时期也正是 \textbf{AI 寒冬}的开始：AI 各个子领域（视觉、专家系统、机器人）都未能兑现早期的宏大承诺，研究经费大幅缩减。

\subsection{1990s--2000s：认知科学的启示与数据驱动的开端}

\subsubsection{认知神经科学的北极星}

尽管处于 AI 寒冬之中，认知和神经科学的研究仍在蓬勃发展，并为计算机视觉指明了方向：

\begin{figure}[H]
\centering
\includegraphics[width=0.95\textwidth]{slides-images/slide-030.jpg}
\caption{视觉识别是视觉智能的基本任务：左侧为自然图像中的视觉识别，右侧展示人脑中专门负责识别面孔、场所、身体部位的特化区域（来自 MIT 的 fMRI 研究）\protect\footnotemark}
\end{figure}
\footnotetext{Slides 第31页（Part 1）。}

\begin{importantbox}{认知科学的三条启示}
\begin{enumerate}
    \item \textbf{场景上下文影响物体识别}：将图像打乱后，即使目标物体的位置和像素完全一致，人类识别该物体的能力也会显著下降（Biederman 实验）
    \item \textbf{视觉处理极其快速}：人类在仅 150 毫秒内就能产生分类信号（Thorpe 实验），这在生物神经网络中仅相当于几个``跳''的处理时间
    \item \textbf{专化的脑区}：人脑发展出了专门识别面孔（FFA）、场所（PPA）和身体部位的特化区域（MIT 1990年代发现）
\end{enumerate}
这些发现告诉我们：物体识别是视觉智能的核心能力，也应该成为计算机视觉的首要研究目标。
\end{importantbox}

\subsubsection{从特征工程到数据集}

1990年代至2000年代初期，计算机视觉领域出现了一系列重要进展：

\begin{itemize}
    \item \textbf{1997年}：Normalized Cuts（图像分割）
    \item \textbf{1999年}：SIFT 特征（尺度不变特征变换）——手工设计的特征提取方法
    \item \textbf{2001年}：Viola-Jones 人脸检测算法——5年后被应用于消费级数码相机的自动对焦功能
    \item \textbf{2004--2007年}：Caltech-101、PASCAL VOC 等小型数据集出现
\end{itemize}

互联网的兴起和数字相机的普及为计算机视觉提供了前所未有的数据资源。这标志着从``靠算法''到``靠数据+算法''的范式转变的开始。

\subsection{本章小结}

计算机视觉的发展经历了从神经科学启发（1950s）到理论框架建立（1970s）、AI 寒冬中的坚持（1980s--1990s），再到数据驱动范式萌芽（2000s）的漫长历程。每一个阶段都为后来的深度学习革命奠定了不可或缺的基础。

%% ========== 深度学习的崛起 ==========

